\subsection*{Вариант 1 для  Рыльской Эвилины Ильгамовны. Дроби \\ \today   \currenttime} 
  \begin{enumerate}
\item Найди число, если его $\frac{1}{4}$  часть равна 32. Нарисуй схему.  \\ \nopagebreak \begin{tikzpicture}
 \draw[step=0.7cm ,gray,very thin] (0,0) grid (18,6);
 % Обводка внешнего контура  \draw[black, thick] (0,0) rectangle (18,6);
 \end{tikzpicture}
\item Найди $\frac{1}{6}$ часть от 84. Нарисуй схему. \\ \nopagebreak \begin{tikzpicture}
 \draw[step=0.7cm ,gray,very thin] (0,0) grid (18,6);
 % Обводка внешнего контура  \draw[black, thick] (0,0) rectangle (18,6);
 \end{tikzpicture}
\item Найди число, если его $\frac{8}{15}$  часть равна 120. Нарисуй схему.  \\ \nopagebreak \begin{tikzpicture}
 \draw[step=0.7cm ,gray,very thin] (0,0) grid (18,6);
 % Обводка внешнего контура  \draw[black, thick] (0,0) rectangle (18,6);
 \end{tikzpicture}
\item Найди $\frac{12}{22}$ часть от 264.  Нарисуй схему. \\ \nopagebreak \begin{tikzpicture}
 \draw[step=0.7cm ,gray,very thin] (0,0) grid (18,6);
 % Обводка внешнего контура  \draw[black, thick] (0,0) rectangle (18,6);
 \end{tikzpicture}
\item Найди число, если его  12 \%  равны 13200. Нарисуй схему.  \\ \nopagebreak \begin{tikzpicture}
 \draw[step=0.7cm ,gray,very thin] (0,0) grid (18,6);
 % Обводка внешнего контура  \draw[black, thick] (0,0) rectangle (18,6);
 \end{tikzpicture}
\item Найди 5 \%  от 8000.  Нарисуй схему. \\ \nopagebreak \begin{tikzpicture}
 \draw[step=0.7cm ,gray,very thin] (0,0) grid (18,6);
 % Обводка внешнего контура  \draw[black, thick] (0,0) rectangle (18,6);
 \end{tikzpicture}
\end{enumerate}  
\newpage
\subsection*{Вариант 2 для  Абелюк Марии. Дроби \\ \today   \currenttime} 
  \begin{enumerate}
\item Найди число, если его $\frac{1}{7}$  часть равна 63. Нарисуй схему.  \\ \nopagebreak \begin{tikzpicture}
 \draw[step=0.7cm ,gray,very thin] (0,0) grid (18,6);
 % Обводка внешнего контура  \draw[black, thick] (0,0) rectangle (18,6);
 \end{tikzpicture}
\item Найди $\frac{1}{9}$ часть от 63. Нарисуй схему. \\ \nopagebreak \begin{tikzpicture}
 \draw[step=0.7cm ,gray,very thin] (0,0) grid (18,6);
 % Обводка внешнего контура  \draw[black, thick] (0,0) rectangle (18,6);
 \end{tikzpicture}
\item Найди число, если его $\frac{2}{11}$  часть равна 22. Нарисуй схему.  \\ \nopagebreak \begin{tikzpicture}
 \draw[step=0.7cm ,gray,very thin] (0,0) grid (18,6);
 % Обводка внешнего контура  \draw[black, thick] (0,0) rectangle (18,6);
 \end{tikzpicture}
\item Найди $\frac{5}{14}$ часть от 70.  Нарисуй схему. \\ \nopagebreak \begin{tikzpicture}
 \draw[step=0.7cm ,gray,very thin] (0,0) grid (18,6);
 % Обводка внешнего контура  \draw[black, thick] (0,0) rectangle (18,6);
 \end{tikzpicture}
\item Найди число, если его  15 \%  равны 7500. Нарисуй схему.  \\ \nopagebreak \begin{tikzpicture}
 \draw[step=0.7cm ,gray,very thin] (0,0) grid (18,6);
 % Обводка внешнего контура  \draw[black, thick] (0,0) rectangle (18,6);
 \end{tikzpicture}
\item Найди 7 \%  от 5600.  Нарисуй схему. \\ \nopagebreak \begin{tikzpicture}
 \draw[step=0.7cm ,gray,very thin] (0,0) grid (18,6);
 % Обводка внешнего контура  \draw[black, thick] (0,0) rectangle (18,6);
 \end{tikzpicture}
\end{enumerate}  
\newpage
\subsection*{Вариант 3 для  Гондарева Ильи. Дроби \\ \today   \currenttime} 
  \begin{enumerate}
\item Найди число, если его $\frac{1}{9}$  часть равна 45. Нарисуй схему.  \\ \nopagebreak \begin{tikzpicture}
 \draw[step=0.7cm ,gray,very thin] (0,0) grid (18,6);
 % Обводка внешнего контура  \draw[black, thick] (0,0) rectangle (18,6);
 \end{tikzpicture}
\item Найди $\frac{1}{2}$ часть от 26. Нарисуй схему. \\ \nopagebreak \begin{tikzpicture}
 \draw[step=0.7cm ,gray,very thin] (0,0) grid (18,6);
 % Обводка внешнего контура  \draw[black, thick] (0,0) rectangle (18,6);
 \end{tikzpicture}
\item Найди число, если его $\frac{6}{9}$  часть равна 54. Нарисуй схему.  \\ \nopagebreak \begin{tikzpicture}
 \draw[step=0.7cm ,gray,very thin] (0,0) grid (18,6);
 % Обводка внешнего контура  \draw[black, thick] (0,0) rectangle (18,6);
 \end{tikzpicture}
\item Найди $\frac{6}{11}$ часть от 66.  Нарисуй схему. \\ \nopagebreak \begin{tikzpicture}
 \draw[step=0.7cm ,gray,very thin] (0,0) grid (18,6);
 % Обводка внешнего контура  \draw[black, thick] (0,0) rectangle (18,6);
 \end{tikzpicture}
\item Найди число, если его  12 \%  равны 6000. Нарисуй схему.  \\ \nopagebreak \begin{tikzpicture}
 \draw[step=0.7cm ,gray,very thin] (0,0) grid (18,6);
 % Обводка внешнего контура  \draw[black, thick] (0,0) rectangle (18,6);
 \end{tikzpicture}
\item Найди 12 \%  от 8400.  Нарисуй схему. \\ \nopagebreak \begin{tikzpicture}
 \draw[step=0.7cm ,gray,very thin] (0,0) grid (18,6);
 % Обводка внешнего контура  \draw[black, thick] (0,0) rectangle (18,6);
 \end{tikzpicture}
\end{enumerate}  
\newpage
\subsection*{Вариант 4 для  Аношина Матвея Антоновича. Дроби \\ \today   \currenttime} 
  \begin{enumerate}
\item Найди число, если его $\frac{1}{10}$  часть равна 40. Нарисуй схему.  \\ \nopagebreak \begin{tikzpicture}
 \draw[step=0.7cm ,gray,very thin] (0,0) grid (18,6);
 % Обводка внешнего контура  \draw[black, thick] (0,0) rectangle (18,6);
 \end{tikzpicture}
\item Найди $\frac{1}{5}$ часть от 55. Нарисуй схему. \\ \nopagebreak \begin{tikzpicture}
 \draw[step=0.7cm ,gray,very thin] (0,0) grid (18,6);
 % Обводка внешнего контура  \draw[black, thick] (0,0) rectangle (18,6);
 \end{tikzpicture}
\item Найди число, если его $\frac{2}{14}$  часть равна 28. Нарисуй схему.  \\ \nopagebreak \begin{tikzpicture}
 \draw[step=0.7cm ,gray,very thin] (0,0) grid (18,6);
 % Обводка внешнего контура  \draw[black, thick] (0,0) rectangle (18,6);
 \end{tikzpicture}
\item Найди $\frac{10}{18}$ часть от 180.  Нарисуй схему. \\ \nopagebreak \begin{tikzpicture}
 \draw[step=0.7cm ,gray,very thin] (0,0) grid (18,6);
 % Обводка внешнего контура  \draw[black, thick] (0,0) rectangle (18,6);
 \end{tikzpicture}
\item Найди число, если его  14 \%  равны 25200. Нарисуй схему.  \\ \nopagebreak \begin{tikzpicture}
 \draw[step=0.7cm ,gray,very thin] (0,0) grid (18,6);
 % Обводка внешнего контура  \draw[black, thick] (0,0) rectangle (18,6);
 \end{tikzpicture}
\item Найди 5 \%  от 3500.  Нарисуй схему. \\ \nopagebreak \begin{tikzpicture}
 \draw[step=0.7cm ,gray,very thin] (0,0) grid (18,6);
 % Обводка внешнего контура  \draw[black, thick] (0,0) rectangle (18,6);
 \end{tikzpicture}
\end{enumerate}  
\newpage
\subsection*{Вариант 5 для  Егоровой Мирославы Даниловны. Дроби \\ \today   \currenttime} 
  \begin{enumerate}
\item Найди число, если его $\frac{1}{4}$  часть равна 8. Нарисуй схему.  \\ \nopagebreak \begin{tikzpicture}
 \draw[step=0.7cm ,gray,very thin] (0,0) grid (18,6);
 % Обводка внешнего контура  \draw[black, thick] (0,0) rectangle (18,6);
 \end{tikzpicture}
\item Найди $\frac{1}{15}$ часть от 45. Нарисуй схему. \\ \nopagebreak \begin{tikzpicture}
 \draw[step=0.7cm ,gray,very thin] (0,0) grid (18,6);
 % Обводка внешнего контура  \draw[black, thick] (0,0) rectangle (18,6);
 \end{tikzpicture}
\item Найди число, если его $\frac{6}{9}$  часть равна 54. Нарисуй схему.  \\ \nopagebreak \begin{tikzpicture}
 \draw[step=0.7cm ,gray,very thin] (0,0) grid (18,6);
 % Обводка внешнего контура  \draw[black, thick] (0,0) rectangle (18,6);
 \end{tikzpicture}
\item Найди $\frac{12}{18}$ часть от 216.  Нарисуй схему. \\ \nopagebreak \begin{tikzpicture}
 \draw[step=0.7cm ,gray,very thin] (0,0) grid (18,6);
 % Обводка внешнего контура  \draw[black, thick] (0,0) rectangle (18,6);
 \end{tikzpicture}
\item Найди число, если его  17 \%  равны 35700. Нарисуй схему.  \\ \nopagebreak \begin{tikzpicture}
 \draw[step=0.7cm ,gray,very thin] (0,0) grid (18,6);
 % Обводка внешнего контура  \draw[black, thick] (0,0) rectangle (18,6);
 \end{tikzpicture}
\item Найди 4 \%  от 8800.  Нарисуй схему. \\ \nopagebreak \begin{tikzpicture}
 \draw[step=0.7cm ,gray,very thin] (0,0) grid (18,6);
 % Обводка внешнего контура  \draw[black, thick] (0,0) rectangle (18,6);
 \end{tikzpicture}
\end{enumerate}  
\newpage
\subsection*{Вариант 6 для  Карнауха Тимофея Павловича. Дроби \\ \today   \currenttime} 
  \begin{enumerate}
\item Найди число, если его $\frac{1}{9}$  часть равна 63. Нарисуй схему.  \\ \nopagebreak \begin{tikzpicture}
 \draw[step=0.7cm ,gray,very thin] (0,0) grid (18,6);
 % Обводка внешнего контура  \draw[black, thick] (0,0) rectangle (18,6);
 \end{tikzpicture}
\item Найди $\frac{1}{3}$ часть от 33. Нарисуй схему. \\ \nopagebreak \begin{tikzpicture}
 \draw[step=0.7cm ,gray,very thin] (0,0) grid (18,6);
 % Обводка внешнего контура  \draw[black, thick] (0,0) rectangle (18,6);
 \end{tikzpicture}
\item Найди число, если его $\frac{2}{12}$  часть равна 24. Нарисуй схему.  \\ \nopagebreak \begin{tikzpicture}
 \draw[step=0.7cm ,gray,very thin] (0,0) grid (18,6);
 % Обводка внешнего контура  \draw[black, thick] (0,0) rectangle (18,6);
 \end{tikzpicture}
\item Найди $\frac{6}{16}$ часть от 96.  Нарисуй схему. \\ \nopagebreak \begin{tikzpicture}
 \draw[step=0.7cm ,gray,very thin] (0,0) grid (18,6);
 % Обводка внешнего контура  \draw[black, thick] (0,0) rectangle (18,6);
 \end{tikzpicture}
\item Найди число, если его  16 \%  равны 16000. Нарисуй схему.  \\ \nopagebreak \begin{tikzpicture}
 \draw[step=0.7cm ,gray,very thin] (0,0) grid (18,6);
 % Обводка внешнего контура  \draw[black, thick] (0,0) rectangle (18,6);
 \end{tikzpicture}
\item Найди 10 \%  от 21000.  Нарисуй схему. \\ \nopagebreak \begin{tikzpicture}
 \draw[step=0.7cm ,gray,very thin] (0,0) grid (18,6);
 % Обводка внешнего контура  \draw[black, thick] (0,0) rectangle (18,6);
 \end{tikzpicture}
\end{enumerate}  
\newpage
\subsection*{Вариант 7 для  Костенко Каролины Романовны. Дроби \\ \today   \currenttime} 
  \begin{enumerate}
\item Найди число, если его $\frac{1}{10}$  часть равна 70. Нарисуй схему.  \\ \nopagebreak \begin{tikzpicture}
 \draw[step=0.7cm ,gray,very thin] (0,0) grid (18,6);
 % Обводка внешнего контура  \draw[black, thick] (0,0) rectangle (18,6);
 \end{tikzpicture}
\item Найди $\frac{1}{9}$ часть от 45. Нарисуй схему. \\ \nopagebreak \begin{tikzpicture}
 \draw[step=0.7cm ,gray,very thin] (0,0) grid (18,6);
 % Обводка внешнего контура  \draw[black, thick] (0,0) rectangle (18,6);
 \end{tikzpicture}
\item Найди число, если его $\frac{4}{8}$  часть равна 32. Нарисуй схему.  \\ \nopagebreak \begin{tikzpicture}
 \draw[step=0.7cm ,gray,very thin] (0,0) grid (18,6);
 % Обводка внешнего контура  \draw[black, thick] (0,0) rectangle (18,6);
 \end{tikzpicture}
\item Найди $\frac{13}{21}$ часть от 273.  Нарисуй схему. \\ \nopagebreak \begin{tikzpicture}
 \draw[step=0.7cm ,gray,very thin] (0,0) grid (18,6);
 % Обводка внешнего контура  \draw[black, thick] (0,0) rectangle (18,6);
 \end{tikzpicture}
\item Найди число, если его  6 \%  равны 6600. Нарисуй схему.  \\ \nopagebreak \begin{tikzpicture}
 \draw[step=0.7cm ,gray,very thin] (0,0) grid (18,6);
 % Обводка внешнего контура  \draw[black, thick] (0,0) rectangle (18,6);
 \end{tikzpicture}
\item Найди 3 \%  от 6000.  Нарисуй схему. \\ \nopagebreak \begin{tikzpicture}
 \draw[step=0.7cm ,gray,very thin] (0,0) grid (18,6);
 % Обводка внешнего контура  \draw[black, thick] (0,0) rectangle (18,6);
 \end{tikzpicture}
\end{enumerate}  
\newpage
\subsection*{Вариант 8 для  Минца Марка Евгеньевича. Дроби \\ \today   \currenttime} 
  \begin{enumerate}
\item Найди число, если его $\frac{1}{9}$  часть равна 45. Нарисуй схему.  \\ \nopagebreak \begin{tikzpicture}
 \draw[step=0.7cm ,gray,very thin] (0,0) grid (18,6);
 % Обводка внешнего контура  \draw[black, thick] (0,0) rectangle (18,6);
 \end{tikzpicture}
\item Найди $\frac{1}{12}$ часть от 156. Нарисуй схему. \\ \nopagebreak \begin{tikzpicture}
 \draw[step=0.7cm ,gray,very thin] (0,0) grid (18,6);
 % Обводка внешнего контура  \draw[black, thick] (0,0) rectangle (18,6);
 \end{tikzpicture}
\item Найди число, если его $\frac{7}{18}$  часть равна 126. Нарисуй схему.  \\ \nopagebreak \begin{tikzpicture}
 \draw[step=0.7cm ,gray,very thin] (0,0) grid (18,6);
 % Обводка внешнего контура  \draw[black, thick] (0,0) rectangle (18,6);
 \end{tikzpicture}
\item Найди $\frac{8}{18}$ часть от 144.  Нарисуй схему. \\ \nopagebreak \begin{tikzpicture}
 \draw[step=0.7cm ,gray,very thin] (0,0) grid (18,6);
 % Обводка внешнего контура  \draw[black, thick] (0,0) rectangle (18,6);
 \end{tikzpicture}
\item Найди число, если его  6 \%  равны 12000. Нарисуй схему.  \\ \nopagebreak \begin{tikzpicture}
 \draw[step=0.7cm ,gray,very thin] (0,0) grid (18,6);
 % Обводка внешнего контура  \draw[black, thick] (0,0) rectangle (18,6);
 \end{tikzpicture}
\item Найди 7 \%  от 11200.  Нарисуй схему. \\ \nopagebreak \begin{tikzpicture}
 \draw[step=0.7cm ,gray,very thin] (0,0) grid (18,6);
 % Обводка внешнего контура  \draw[black, thick] (0,0) rectangle (18,6);
 \end{tikzpicture}
\end{enumerate}  
\newpage
\subsection*{Вариант 9 для  Полякова Никиты. Дроби \\ \today   \currenttime} 
  \begin{enumerate}
\item Найди число, если его $\frac{1}{7}$  часть равна 21. Нарисуй схему.  \\ \nopagebreak \begin{tikzpicture}
 \draw[step=0.7cm ,gray,very thin] (0,0) grid (18,6);
 % Обводка внешнего контура  \draw[black, thick] (0,0) rectangle (18,6);
 \end{tikzpicture}
\item Найди $\frac{1}{5}$ часть от 55. Нарисуй схему. \\ \nopagebreak \begin{tikzpicture}
 \draw[step=0.7cm ,gray,very thin] (0,0) grid (18,6);
 % Обводка внешнего контура  \draw[black, thick] (0,0) rectangle (18,6);
 \end{tikzpicture}
\item Найди число, если его $\frac{8}{18}$  часть равна 144. Нарисуй схему.  \\ \nopagebreak \begin{tikzpicture}
 \draw[step=0.7cm ,gray,very thin] (0,0) grid (18,6);
 % Обводка внешнего контура  \draw[black, thick] (0,0) rectangle (18,6);
 \end{tikzpicture}
\item Найди $\frac{11}{18}$ часть от 198.  Нарисуй схему. \\ \nopagebreak \begin{tikzpicture}
 \draw[step=0.7cm ,gray,very thin] (0,0) grid (18,6);
 % Обводка внешнего контура  \draw[black, thick] (0,0) rectangle (18,6);
 \end{tikzpicture}
\item Найди число, если его  13 \%  равны 27300. Нарисуй схему.  \\ \nopagebreak \begin{tikzpicture}
 \draw[step=0.7cm ,gray,very thin] (0,0) grid (18,6);
 % Обводка внешнего контура  \draw[black, thick] (0,0) rectangle (18,6);
 \end{tikzpicture}
\item Найди 4 \%  от 3200.  Нарисуй схему. \\ \nopagebreak \begin{tikzpicture}
 \draw[step=0.7cm ,gray,very thin] (0,0) grid (18,6);
 % Обводка внешнего контура  \draw[black, thick] (0,0) rectangle (18,6);
 \end{tikzpicture}
\end{enumerate}  
\newpage
\subsection*{Вариант 10 для  Саргсян Эмилии Арамовны. Дроби \\ \today   \currenttime} 
  \begin{enumerate}
\item Найди число, если его $\frac{1}{10}$  часть равна 30. Нарисуй схему.  \\ \nopagebreak \begin{tikzpicture}
 \draw[step=0.7cm ,gray,very thin] (0,0) grid (18,6);
 % Обводка внешнего контура  \draw[black, thick] (0,0) rectangle (18,6);
 \end{tikzpicture}
\item Найди $\frac{1}{3}$ часть от 24. Нарисуй схему. \\ \nopagebreak \begin{tikzpicture}
 \draw[step=0.7cm ,gray,very thin] (0,0) grid (18,6);
 % Обводка внешнего контура  \draw[black, thick] (0,0) rectangle (18,6);
 \end{tikzpicture}
\item Найди число, если его $\frac{9}{13}$  часть равна 117. Нарисуй схему.  \\ \nopagebreak \begin{tikzpicture}
 \draw[step=0.7cm ,gray,very thin] (0,0) grid (18,6);
 % Обводка внешнего контура  \draw[black, thick] (0,0) rectangle (18,6);
 \end{tikzpicture}
\item Найди $\frac{8}{16}$ часть от 128.  Нарисуй схему. \\ \nopagebreak \begin{tikzpicture}
 \draw[step=0.7cm ,gray,very thin] (0,0) grid (18,6);
 % Обводка внешнего контура  \draw[black, thick] (0,0) rectangle (18,6);
 \end{tikzpicture}
\item Найди число, если его  4 \%  равны 3600. Нарисуй схему.  \\ \nopagebreak \begin{tikzpicture}
 \draw[step=0.7cm ,gray,very thin] (0,0) grid (18,6);
 % Обводка внешнего контура  \draw[black, thick] (0,0) rectangle (18,6);
 \end{tikzpicture}
\item Найди 12 \%  от 27600.  Нарисуй схему. \\ \nopagebreak \begin{tikzpicture}
 \draw[step=0.7cm ,gray,very thin] (0,0) grid (18,6);
 % Обводка внешнего контура  \draw[black, thick] (0,0) rectangle (18,6);
 \end{tikzpicture}
\end{enumerate}  
\newpage
\subsection*{Вариант 11 для  Строкатовой Киры Павловны. Дроби \\ \today   \currenttime} 
  \begin{enumerate}
\item Найди число, если его $\frac{1}{7}$  часть равна 70. Нарисуй схему.  \\ \nopagebreak \begin{tikzpicture}
 \draw[step=0.7cm ,gray,very thin] (0,0) grid (18,6);
 % Обводка внешнего контура  \draw[black, thick] (0,0) rectangle (18,6);
 \end{tikzpicture}
\item Найди $\frac{1}{10}$ часть от 70. Нарисуй схему. \\ \nopagebreak \begin{tikzpicture}
 \draw[step=0.7cm ,gray,very thin] (0,0) grid (18,6);
 % Обводка внешнего контура  \draw[black, thick] (0,0) rectangle (18,6);
 \end{tikzpicture}
\item Найди число, если его $\frac{3}{6}$  часть равна 18. Нарисуй схему.  \\ \nopagebreak \begin{tikzpicture}
 \draw[step=0.7cm ,gray,very thin] (0,0) grid (18,6);
 % Обводка внешнего контура  \draw[black, thick] (0,0) rectangle (18,6);
 \end{tikzpicture}
\item Найди $\frac{8}{16}$ часть от 128.  Нарисуй схему. \\ \nopagebreak \begin{tikzpicture}
 \draw[step=0.7cm ,gray,very thin] (0,0) grid (18,6);
 % Обводка внешнего контура  \draw[black, thick] (0,0) rectangle (18,6);
 \end{tikzpicture}
\item Найди число, если его  6 \%  равны 10200. Нарисуй схему.  \\ \nopagebreak \begin{tikzpicture}
 \draw[step=0.7cm ,gray,very thin] (0,0) grid (18,6);
 % Обводка внешнего контура  \draw[black, thick] (0,0) rectangle (18,6);
 \end{tikzpicture}
\item Найди 8 \%  от 9600.  Нарисуй схему. \\ \nopagebreak \begin{tikzpicture}
 \draw[step=0.7cm ,gray,very thin] (0,0) grid (18,6);
 % Обводка внешнего контура  \draw[black, thick] (0,0) rectangle (18,6);
 \end{tikzpicture}
\end{enumerate}  
\newpage
\subsection*{Вариант 12 для  Тавриной Марии Ивановны. Дроби \\ \today   \currenttime} 
  \begin{enumerate}
\item Найди число, если его $\frac{1}{11}$  часть равна 44. Нарисуй схему.  \\ \nopagebreak \begin{tikzpicture}
 \draw[step=0.7cm ,gray,very thin] (0,0) grid (18,6);
 % Обводка внешнего контура  \draw[black, thick] (0,0) rectangle (18,6);
 \end{tikzpicture}
\item Найди $\frac{1}{3}$ часть от 18. Нарисуй схему. \\ \nopagebreak \begin{tikzpicture}
 \draw[step=0.7cm ,gray,very thin] (0,0) grid (18,6);
 % Обводка внешнего контура  \draw[black, thick] (0,0) rectangle (18,6);
 \end{tikzpicture}
\item Найди число, если его $\frac{6}{16}$  часть равна 96. Нарисуй схему.  \\ \nopagebreak \begin{tikzpicture}
 \draw[step=0.7cm ,gray,very thin] (0,0) grid (18,6);
 % Обводка внешнего контура  \draw[black, thick] (0,0) rectangle (18,6);
 \end{tikzpicture}
\item Найди $\frac{13}{22}$ часть от 286.  Нарисуй схему. \\ \nopagebreak \begin{tikzpicture}
 \draw[step=0.7cm ,gray,very thin] (0,0) grid (18,6);
 % Обводка внешнего контура  \draw[black, thick] (0,0) rectangle (18,6);
 \end{tikzpicture}
\item Найди число, если его  11 \%  равны 24200. Нарисуй схему.  \\ \nopagebreak \begin{tikzpicture}
 \draw[step=0.7cm ,gray,very thin] (0,0) grid (18,6);
 % Обводка внешнего контура  \draw[black, thick] (0,0) rectangle (18,6);
 \end{tikzpicture}
\item Найди 11 \%  от 18700.  Нарисуй схему. \\ \nopagebreak \begin{tikzpicture}
 \draw[step=0.7cm ,gray,very thin] (0,0) grid (18,6);
 % Обводка внешнего контура  \draw[black, thick] (0,0) rectangle (18,6);
 \end{tikzpicture}
\end{enumerate}  
\newpage
